% GENERATED FILE. Edit canonical lessons, then run npm run generate:books.
% canonical-source-hash: fnv1a64:868983b6a640d151
% canonical-lessons: TA-C313-itarkitaiyil, TA-C313-utanatiyaka, TA-C313-patippatiyaka, TA-C313-oruvaliyaka, TA-C313-munnataka

\chapter{Time Words for a Story}
\label{ch:ta-time-words-for-a-story}
\hlchaptermodality{313}

\begin{chapteropening}
\textbf{By the end of this chapter:} \emph{I can link the events of a story in time.}

Complete the last lesson of chapter 313: I can link the events of a story in time.
\end{chapteropening}

\section[\texorpdfstring{\ta{இதற்கிடையில்} (itaṟkiṭaiyil)}{இதற்கிடையில் (itaṟkiṭaiyil)}]{\ta{இதற்கிடையில்} (itaṟkiṭaiyil) — meanwhile}

\label{lesson:TA-C313-itarkitaiyil}



\begin{quote}
Before the new one: say the Tamil for to recite, then the Tamil for oral tradition.
\end{quote}

\subsection*{You'll want to know: \ta{இதற்கிடையில்}}
\textbf{\ta{இதற்கிடையில்}} — \emph{itaṟkiṭaiyil} — \textquotedblleft{}meanwhile\textquotedblright{}.

\begin{quote}
\ta{இதற்கிடையில்}, \ta{காவலர்கள்} \ta{வீட்டுக்கு} \ta{வந்தார்கள்}.
\end{quote}

\begin{quote}
\emph{itaṟkiṭaiyil, kāvalarkaḷ vīṭṭukku vantārkaḷ.}
\end{quote}

\begin{quote}
Meanwhile, the police came to the house.
\end{quote}

\ta{இது} + \ta{இடை}, the space between, + -\ta{இல்}: in the meantime.

\begin{cousinweb}
\textbf{\ta{இதற்கிடையில்}} (\emph{itaṟkiṭaiyil}) is a Tamil phrase grown into one word, built from \ta{இதற்கு}, \textquotedblleft{}to this\textquotedblright{}, and \ta{இடையில்}, \textquotedblleft{}in between\textquotedblright{}, from native \ta{இடை}, \textquotedblleft{}space, interval\textquotedblright{} (DEDR 434).
\end{cousinweb}

\begin{grammarlens}[title={What you've built}]
One more word for telling what happened.
\end{grammarlens}

\subsection*{Your turn}
\begin{itemize}
\raggedright
  \item \emph{Say it:} \emph{itaṟkiṭaiyil}
  \item \emph{Say it:} \emph{itaṟkiṭaiyil}, once more
  \item \emph{Say it:} \emph{itaṟkiṭaiyil, kāvalarkaḷ vīṭṭukku vantārkaḷ.} — Meanwhile, the police came to the house.
  \item \emph{Recall:} say the Tamil for to recite, then the Tamil for oral tradition, then say \emph{itaṟkiṭaiyil} again
\end{itemize}

\begin{tcolorbox}[breakable,colback=teal!4,colframe=teal!35!black,title={Before you move on}]
What does \ta{இதற்கிடையில்} mean? (\textquotedblleft{}Meanwhile\textquotedblright{}.) Say it once more.
\end{tcolorbox}
\section[\texorpdfstring{\ta{உடனடியாக} (uṭaṉaṭiyāka)}{உடனடியாக (uṭaṉaṭiyāka)}]{\ta{உடனடியாக} (uṭaṉaṭiyāka) — immediately}

\label{lesson:TA-C313-utanatiyaka}



\begin{quote}
Before the new one: say the Tamil for oral tradition, then the Tamil for meanwhile.
\end{quote}

\subsection*{You'll want to know: \ta{உடனடியாக}}
\textbf{\ta{உடனடியாக}} — \emph{uṭaṉaṭiyāka} — \textquotedblleft{}immediately\textquotedblright{}.

\begin{quote}
\ta{தீயணைப்பு} \ta{வீரர்கள்} \ta{உடனடியாக} \ta{வந்தார்கள்}.
\end{quote}

\begin{quote}
\emph{tīyaṇaippu vīrarkaḷ uṭaṉaṭiyāka vantārkaḷ.}
\end{quote}

\begin{quote}
The firefighters came immediately.
\end{quote}

Stronger and more formal than \ta{உடனே}.

\begin{cousinweb}
\textbf{\ta{உடனடியாக}} (\emph{uṭaṉaṭiyāka}) is built on native \ta{உடன்}, \textquotedblleft{}together with, at once\textquotedblright{} (DEDR 945), with \ta{அடி} and the adverb ending -\ta{ஆக}.
\end{cousinweb}

\begin{grammarlens}[title={What you've built}]
One more word for telling what happened.
\end{grammarlens}

\subsection*{Your turn}
\begin{itemize}
\raggedright
  \item \emph{Say it:} \emph{uṭaṉaṭiyāka}
  \item \emph{Say it:} \emph{uṭaṉaṭiyāka}, once more
  \item \emph{Say it:} \emph{tīyaṇaippu vīrarkaḷ uṭaṉaṭiyāka vantārkaḷ.} — The firefighters came immediately.
  \item \emph{Recall:} say the Tamil for oral tradition, then the Tamil for meanwhile, then say \emph{uṭaṉaṭiyāka} again
\end{itemize}

\begin{tcolorbox}[breakable,colback=teal!4,colframe=teal!35!black,title={Before you move on}]
What does \ta{உடனடியாக} mean? (\textquotedblleft{}Immediately\textquotedblright{}.) Say it once more.
\end{tcolorbox}
\section[\texorpdfstring{\ta{படிப்படியாக} (paṭippaṭiyāka)}{படிப்படியாக (paṭippaṭiyāka)}]{\ta{படிப்படியாக} (paṭippaṭiyāka) — gradually; step by step}

\label{lesson:TA-C313-patippatiyaka}



\begin{quote}
Before the new one: say the Tamil for meanwhile, then the Tamil for immediately.
\end{quote}

\subsection*{You'll want to know: \ta{படிப்படியாக}}
\textbf{\ta{படிப்படியாக}} — \emph{paṭippaṭiyāka} — \textquotedblleft{}gradually; step by step\textquotedblright{}.

\begin{quote}
\ta{மழை} \ta{படிப்படியாகக்} \ta{குறைந்தது}.
\end{quote}

\begin{quote}
\emph{maḻai paṭippaṭiyākak kuṟaintatu.}
\end{quote}

\begin{quote}
The rain gradually eased off.
\end{quote}

\ta{படி}, step, doubled: step by step.

\begin{cousinweb}
\textbf{\ta{படிப்படியாக}} (\emph{paṭippaṭiyāka}) doubles native \ta{படி}, \textquotedblleft{}a step, a grade\textquotedblright{} (DEDR 3850), with the adverb ending -\ta{ஆக}.
\end{cousinweb}

\begin{grammarlens}[title={What you've built}]
One more word for telling what happened.
\end{grammarlens}

\subsection*{Your turn}
\begin{itemize}
\raggedright
  \item \emph{Say it:} \emph{paṭippaṭiyāka}
  \item \emph{Say it:} \emph{paṭippaṭiyāka}, once more
  \item \emph{Say it:} \emph{maḻai paṭippaṭiyākak kuṟaintatu.} — The rain gradually eased off.
  \item \emph{Recall:} say the Tamil for meanwhile, then the Tamil for immediately, then say \emph{paṭippaṭiyāka} again
\end{itemize}

\begin{tcolorbox}[breakable,colback=teal!4,colframe=teal!35!black,title={Before you move on}]
What does \ta{படிப்படியாக} mean? (\textquotedblleft{}Gradually; step by step\textquotedblright{}.) Say it once more.
\end{tcolorbox}
\section[\texorpdfstring{\ta{ஒருவழியாக} (oruvaḻiyāka)}{ஒருவழியாக (oruvaḻiyāka)}]{\ta{ஒருவழியாக} (oruvaḻiyāka) — at long last}

\label{lesson:TA-C313-oruvaliyaka}



\begin{quote}
Before the new one: say the Tamil for immediately, then the Tamil for gradually.
\end{quote}

\subsection*{You'll want to know: \ta{ஒருவழியாக}}
\textbf{\ta{ஒருவழியாக}} — \emph{oruvaḻiyāka} — \textquotedblleft{}at long last\textquotedblright{}.

\begin{quote}
\ta{ஒருவழியாக}, \ta{பேருந்து} \ta{வந்தது}.
\end{quote}

\begin{quote}
\emph{oruvaḻiyāka, pēruntu vantatu.}
\end{quote}

\begin{quote}
At long last, the bus came.
\end{quote}

Said with relief after a wait or a struggle.

\begin{cousinweb}
\textbf{\ta{ஒருவழியாக}} (\emph{oruvaḻiyāka}) is built from \ta{ஒரு}, \textquotedblleft{}one\textquotedblright{}, and native \ta{வழி}, \textquotedblleft{}way, means\textquotedblright{} (DEDR 5297), with the adverb ending -\ta{ஆக}: somehow, one way or another.
\end{cousinweb}

\begin{grammarlens}[title={What you've built}]
One more word for telling what happened.
\end{grammarlens}

\subsection*{Your turn}
\begin{itemize}
\raggedright
  \item \emph{Say it:} \emph{oruvaḻiyāka}
  \item \emph{Say it:} \emph{oruvaḻiyāka}, once more
  \item \emph{Say it:} \emph{oruvaḻiyāka, pēruntu vantatu.} — At long last, the bus came.
  \item \emph{Recall:} say the Tamil for immediately, then the Tamil for gradually, then say \emph{oruvaḻiyāka} again
\end{itemize}

\begin{tcolorbox}[breakable,colback=teal!4,colframe=teal!35!black,title={Before you move on}]
What does \ta{ஒருவழியாக} mean? (\textquotedblleft{}At long last\textquotedblright{}.) Say it once more.
\end{tcolorbox}
\section[\texorpdfstring{\ta{முன்னதாக} (muṉṉatāka)}{முன்னதாக (muṉṉatāka)}]{\ta{முன்னதாக} (muṉṉatāka) — earlier; beforehand}

\label{lesson:TA-C313-munnataka}



\begin{quote}
Before the new one: say the Tamil for gradually, then the Tamil for at long last.
\end{quote}

\subsection*{You'll want to know: \ta{முன்னதாக}}
\textbf{\ta{முன்னதாக}} — \emph{muṉṉatāka} — \textquotedblleft{}earlier; beforehand\textquotedblright{}.

\begin{quote}
\ta{விழா} \ta{ஒரு} \ta{மணி} \ta{நேரம்} \ta{முன்னதாகத்} \ta{தொடங்கியது}.
\end{quote}

\begin{quote}
\emph{viḻā oru maṇi nēram muṉṉatākat toṭaṅkiyatu.}
\end{quote}

\begin{quote}
The ceremony started an hour earlier.
\end{quote}

Also beforehand: \ta{முன்னதாகவே} \ta{சொன்னேன்}, I said so in advance.

\begin{cousinweb}
\textbf{\ta{முன்னதாக}} (\emph{muṉṉatāka}) is built on native \ta{முன்}, \textquotedblleft{}in front, before\textquotedblright{} (DEDR 5020), with the adverb ending -\ta{ஆக}.
\end{cousinweb}

\begin{grammarlens}[title={What you've built}]
One more word for telling what happened.
\end{grammarlens}

\subsection*{Your turn}
\begin{itemize}
\raggedright
  \item \emph{Say it:} \emph{muṉṉatāka}
  \item \emph{Say it:} \emph{muṉṉatāka}, once more
  \item \emph{Say it:} \emph{viḻā oru maṇi nēram muṉṉatākat toṭaṅkiyatu.} — The ceremony started an hour earlier.
  \item \emph{Recall:} say the Tamil for gradually, then the Tamil for at long last, then say \emph{muṉṉatāka} again
\end{itemize}

\begin{tcolorbox}[breakable,colback=teal!4,colframe=teal!35!black,title={Before you move on}]
What does \ta{முன்னதாக} mean? (\textquotedblleft{}Earlier; beforehand\textquotedblright{}.) Say it once more.
\end{tcolorbox}
